\documentclass[11pt]{article}
\usepackage{mystyle}

\title{Rosen --- \S 8.6: Applications of Inclusion--Exclusion\\ \large Selected Exercises}
\author{Alston}
\date{Semester 2, 2026}

\begin{document}
\maketitle

% ======================================================================
\begin{exercise}{5}
    Find the number of primes less than 200 using the principle of
    inclusion--exclusion.
\end{exercise}

% ======================================================================
\begin{exercise}{6}
    An integer is called \emph{squarefree} if it is not divisible by
    the square of a positive integer greater than 1. Find the number
    of squarefree positive integers less than 100.
\end{exercise}

% ======================================================================
\begin{exercise}{8}
    How many onto functions are there from a set with seven elements
    to one with five elements?
\end{exercise}

% ======================================================================
\begin{exercise}{14}
    What is the probability that none of 10 people receives the
    correct hat if a hatcheck person hands their hats back randomly?
\end{exercise}

% ======================================================================
\begin{exercise}{17}
    How many ways can the digits 0, 1, 2, 3, 4, 5, 6, 7, 8, 9 be
    arranged so that no even digit is in its original position?
\end{exercise}

% ======================================================================
\begin{exercise}{18}
    Use a combinatorial argument to show that the sequence $\{D_n\}$,
    where $D_n$ denotes the number of derangements of $n$ objects,
    satisfies the recurrence relation
    \[
    D_n = (n-1)(D_{n-1} + D_{n-2})
    \]
    for $n \geq 2$.
    [\emph{Hint}: Note that there are $n-1$ choices for the first
    element $k$ of a derangement. Consider separately the derangements
    that start with $k$ that do and do not have 1 in the $k$th
    position.]
\end{exercise}

% ======================================================================
\begin{exercise}{19}
    Use Exercise 18 to show that
    \[
    D_n = nD_{n-1} + (-1)^n
    \]
    for $n \geq 1$.
\end{exercise}

% ======================================================================
\begin{exercise}{23}
    Use the principle of inclusion--exclusion to derive a formula for
    $\varphi(n)$ when the prime factorization of $n$ is
    $n = p_1^{a_1} p_2^{a_2} \cdots p_m^{a_m}$.
\end{exercise}

% ======================================================================
\begin{exercise}{24}
    Show that if $n$ is a positive integer, then
    \[
    n! = \binom{n}{0} D_n + \binom{n}{1} D_{n-1} + \cdots +
    \binom{n}{n-1} D_1 + \binom{n}{n} D_0,
    \]
    where $D_k$ is the number of derangements of $k$ objects.
\end{exercise}

% ======================================================================
\begin{exercise}{27}
    \textit{Recall (Theorem 1): the number of onto functions from a
    set with $m$ elements to a set with $n$ elements is
    \[
    n^m - \binom{n}{1}(n-1)^m + \binom{n}{2}(n-2)^m - \cdots +
    (-1)^{n-1}\binom{n}{n-1} \cdot 1^m.
    \]}

    Prove Theorem 1.
\end{exercise}

\end{document}